\scanpage{85}

% top, two lines of cancelled opening
\uline{NB} \deleted{La correspondance $C\longmapsto X\setminus\{C\cup\{\unsure{a|\cdot}\}\}$ entre carrés de $X$ et}

% upper and middle, lines 2–9
Appelons \uline{anticarré} l'opposé d'un carré, i.e. un graphe somme de deux segments. Les anticarrés dans $X$ sont donc les carrés dans $X^\circ$, qui est un pavé, donc ils correspondent aux pts de $X^\circ$, en associant à \unsure{tt|$\forall$} $a\in X^\circ$ l'ens des $4$ pts de $X^\circ$ qui sont $X^\circ$-\unsure{non}liés à $a$, i.e. qui sont $X$-liés à $a$. \deleted{Interpréter les} Donc

% middle, lines 10–15
la correspondance $C\longmapsto(X\setminus C\cup\{a=\text{centre de }C\}\unsure{\}})$ entre carrés de $X$ et anticarrés de $X$ est bijective. Si $C$ et $C'$ sont associés ainsi, $C$ et $C'^\circ$ sont des carrés en dualité (les sommets de l'un étant les côtés de l'autre.

% lower, lines 16–22
Il y a une corr. fonctorielle évidente qui associe à chaque carré $C$ un pavé \inserted{$\mathrm{Pav}(C)$} le contenant (son ``enveloppe'') dont l'ens des sommets est $C\amalg C^{*\circ}\amalg\{a\}_{\text{centre}}$. Deux points de $C$ ($C^{*\circ}$) sont liés ssi ils sont liés dans $C$ (\deleted{\unsure{non}}\inserted{resp} liés dans $C^{*\circ}$ i.e non liés dans $C^*$), un pt \inserted{$x$} de $C$ est lié à un pt \inserted{$x^{*}$} de $C^{*\circ}$ ssi \deleted{\unsure{sont}} $x,x^*$ incidents, enfin $a$ est non lié aux pts de $C$, lié aux

% archive sheet number: 84
